% Laget automatisk av elevark/lag_elevark.py fra sim/ballistisk-pendel.html. Endringer her blir overskrevet.
\documentclass[11pt]{article}
\usepackage{elevark}
\begin{document}
\elevarkhode{Fysikk 1 · Bevegelsesmengde og støt}{Ballistisk pendel}{https://lektorbo.no/sim/ballistisk-pendel.html}{Simuleringen}

En kule fra et luftgevær er så rask at den er vanskelig å måle direkte. Skyt den inn i en trekloss som henger i to snorer på 2,0~m. Klossen svinger ut til en høyde $h$. Kan du finne farten til kula ut fra $h$?

\section*{Oppgaver}

Del forløpet i to: selve støtet, og svingningen etterpå. Tenk over hva som er bevart i hver del.

\begin{oppgaveliste}
\oppgave \fase{gjett}{Gjett}Hvor stor andel av den kinetiske energien til kula tror du klossen og kula har rett etter støtet? 100~\%, 50~\%, 10~\% eller 1~\%?
\svar
\oppgave \fase{test}{Test}Skyt en kule på 1,0~g med 200~m/s inn i klossen på 0,20~kg. Les av høyden og andelen kinetisk energi som er igjen. Hvor ble resten av energien av?
\svar
\oppgave \fase{forklar}{Forklar}Hvorfor kan vi ikke bruke energibevaring fra kula før støtet til klossen på toppen? Hvilken bevaringslov gjelder i støtet, og hvilken gjelder i svingningen etterpå?
\svar
\oppgave \fase{test}{Regn}Vis at farten til kula er $v = (m + M)\sqrt{2gh} / m$. Sjekk formelen med tallene fra oppgave 2.
\svar
\oppgave \fase{test}{Mål}Kryss av for «Ukjent fart» og skyt. Les av høyden på målestaven og regn ut farten til kula. Sjekk svaret med «Vis farten». Prøv med flere ukjente kuler.
\svar
\oppgave \fase{forklar}{Forklar}Hva skjer med høyden når du gjør klossen tyngre? Hvorfor kan en for tung kloss gjøre målingen unøyaktig, og hvorfor kan en for lett kloss være et problem?
\svar
\oppgave \fase{utfordring}{Utfordring}Vis at andelen kinetisk energi som er igjen etter støtet, er $m / (m + M)$. Hvorfor avhenger den ikke av farten til kula?
\svar
\end{oppgaveliste}

\end{document}
